\section{"ЛУЧШИЕ УЗЛЫ. ПОДБОРКА ОТ АНАЛИТИКА".- Романенко Е.А..- v3.0, 21.08.2019}

Книга представляет собой подборку небольшого числа узлов, которые, по мнению автора, могут найти наиболее частое применение.

\graphicspath{{\currentpath}}

\subsection{Классификация узлов, предложенная автором}

\begin{enumerate}
\item Соединительные узлы
\begin{enumerate}
\item Легкоразвязываемые соединительные узлы.
Внутри автор выделяет регулируемые узлы
\item Условно-развязываемые соединительные узлы
Часть узлов этой группы автор выделил в "УНИВЕРСАЛЬНЫЕ (любая верёвка/ЛЕСКА)"
\end{enumerate}
\item Петли
\begin{enumerate}
\item Срединные петли
\item Незатягивающиеся петли (концевые)
\begin{itemize}
\item БЫСТРОСБРОСЫ
\item ЛЕГКОРАЗВЯЗЫВАЕМЫЕ НП
\item УСЛОВНО-РАЗВЯЗЫВАЕМЫЕ НП
\item УНИВЕРСАЛЬНЫЕ НП (любая верёвка/ЛЕСКА)
\item ДВОЙНЫЕ НЕЗАТЯГИВАЮЩИЕСЯ ПЕТЛИ
\item Двойные условно-развязываемые НП
\end{itemize}
\item Регулируемые петли
\begin{enumerate}
\item РП с КК-«Опорой»
\item РП с ХК-«Опорой»
\end{enumerate}
\item Петли-защёлки
\item Затягивающиеся петли
\begin{itemize}
\item ДЛЯ ВЕРЁВКИ
\item УНИВЕРСАЛЬНЫЕ (любая верёвка/ЛЕСКА)
\end{itemize}
\item Петли-удавки
\end{enumerate}
\item Крепёжные узлы (вяжутся вокруг Опоры)
Внутри автор выделяет
\begin{itemize}
\item ШТЫКИ
\item ОБЖИМНЫЕ / СТЯГИВАЮЩИЕ
\item УЗЛЫ ДЛЯ КОЛЬЦА («RING HITCHES») – [рин-хичи]
\item СХВАТЫВАЮЩИЕ УЗЛЫ
\item *БУКСИРОВОЧНЫЕ* УЗЛЫ
\end{itemize}
Но список ими не исчерпывается.
\item ОСОБЫЕ УЗЛЫ
\begin{itemize}
\item РЕВЕРСНО-СБРАСЫВАЕМЫЕ УЗЛЫ, или [САМОСБРОСЫ]
\item ПОЛИСПАСТЫ
\item Стопорные узлы, или [узлы для утолщения троса]
\item РАЗНОЕ - сюда автор включил гачные, бочечные, "досочные" и другие узко специализированные узлы
\end{itemize}
Кроме того вне подгрупп идёт амфорный узел.
\end{enumerate}

\paragraph{Пояснения к некоторым группам и взаимосвязи групп}
Легкоразвязываемые узлы - это узлы, которые можно развязать без инструмента.

Условно-развязываемые узлы - это узлы, развязываемость которых без инструмента зависит от ряда условий, таких как свойства верёвки и величина рабочей нагрузкки.

Легкоразвязываемость крепёжных узлов обеспечивается либо конструкцией самого узла, либо бантиком. Автор не считает это классификационным признаком.

Регулируемые узлы, по автору, строются на базе крепёжных узлов, завязанных каждой компонентой вокруг ходового конца другой компоненты. Между узлами идут две верёвки. Т.о. при ослабленных базовых узлах ходовые концы могут двигаться через них, а при затянутых - крепко схвачены. Это позволяет изменять длину регулируемого узла. Логично, что регулируемые узлы должны строиться на основе легкоразвязываемых крепёжных узлов, и как следствие - сами являются легкоразвязываемыми.

В подгруппу универсальные условно-развязываемые соединительные узлы попали узлы, состоящие из двух крепёжных узлов, завязанных вокруг коренных (а не ходовых, как у регулируемых узлов) концов другой компоненты. Но это, очевидно, не классификационный признак. Рыбацкий узел, устроеный аналогично, в эту подгруппу не включён.

Срединными являются узлы, которые можно завязать без доступа к концам верёвки. Т.е. изотопные неузлу.

Срединные петли могут быть только незатягивающимися, причём незатягивающимися относительно обоих концов верёвки. В противном случае, при натяжении верёвки петля исчезнет, а узел выродится в неузел или превратится в другой тип узла. Если перед натяжением петля накинута на опору, то опора становится компонентой, необходимой для сохранения топологии, а, следовательно, такой узел должен быть отнесён к срединным крепёжным узлам, а не к петлям.

При "разрезании" срединной петли - получается соединительный узел.

"Быстросброс" должен легко отдаваться  одним рывком за ХК, не оставляя «пост-узла» + с как можно меньшим захлёстом/перехлёстом («пост-свивкой») концов.
Все [самосбросы], завязанные с коротким ХК, являются [быстросбросами] также.

Быстросбрасываемость всех названных автором быстросбросов обеспечивается бантиком.

Регулируемые петли – это НП изменяемого размера. Относятся к «фиксаторам» («фиксаторным узлам», в данном случае).  Затягивающимися петлями (ЗП) РП быть не могут. Регулировать все РП возможно только в ненагруженном состоянии.

Петли-защёлки – это следующая промежуточная ступень между НП и ЗП.
Между РП и ПЗ, есть принципиальная разница: время, необходимое для изменения размера РП  (особенно класса 1)  –  оч. малó; время, необходимое для изменения размера ПЗ  –  довольно значительно, – счёт идёт уже даже не на секунды – десятки секунд, и львиная доля этого времени затрачивается на существенное ослабление базового узла

По сути, [ПЗ]  –  это [однократно-регулируемые петли класса 1].
А к [однократно-регулируемым петлям класса 2] можно отнести узлы Голландский Булинь, Брусвилис, Шахтёрская петля. – Но не к [ПЗ], т.к. они:
а) защёлкиваются не за ХК, а за КК  (просто приложением рабочей тяги);
б) в незащёлкнутом виде  не могут работать в качестве ЗП  (по причине «а»).

По аналогии, однократно регулируемый узел можно получить, если взять в качестве базового крепёжный узел, не являющийся быстроразвязывающимся.

Затягивающиеся петли одновременно являются универсальными [крепёжными] узлами.

К петлям-удавкам автор относит затягивающиеся петли, которые под рабочей нагрузкой не затрудняют собственное затягивание, а потому максимально легко затягиваются, а при снятии нагрузки - легко ослабляются.

Стягивающими автор называет узлы, которые используются для не для крепления верёвки к опоре, а для связывания деталей опоры вместе.

"Узлы для кольца"  означает:  "Узлы для относительно тонкой Опоры"  –  т.е. для такой Опоры, диаметр которой превышает диаметр верёвки не более чем в 4 раза. В  данной  подгруппе представлены такие хичи, применять которые можно исключительно на тонкой Опоре  (из соображений надёжности), – или же такие, которые  в этом качестве применять предпочтительнее, т.к. именно здесь проявляются их преимущества.

[Схватывающие] – это такой род  хичей, которые без тяги можно свободно передвигать вдоль Опоры, за «тело» узла, – и намертво фиксирующиеся  ("схватывающие" Опору)  при приложении рабочей нагрузки, под тягой.  

*БУКСИРОВОЧНЫЕ* УЗЛЫ - Сохраняют лёгкость развязывания и после экстремально сильного нагружения.

РЕВЕРСНО-СБРАСЫВАЕМЫЕ УЗЛЫ, или [САМОСБРОСЫ]`
Требования, предъявляемые к РСУ.  ...Или  какими  качествами  необходимо обладать *узлу с бантиком*, чтобы претендовать на это высокое звание. – Итак:
1. Самое  важное, я считаю, качество, определяющее надёжность самосброса от несанкционированного сброса – это возможно более "мёртвый" зажим бантика сбросового конца тягой на рабочем (спусковом) конце. Это есть «Критерий №1».
2. На 2-м, по важности, месте находится лёгкость/беспроблемность штатного сброса. – Т.е. в какой мере "мертвый" зажим бантика обратим при реверсе тяги. (...И вот здесь-то и "сыпятся" все остальные узлы-претенденты. – Не проходят: Бегущий Простой, Травяная +Бурлацкая  п-с-б. ...Да и Калмыцкий, в общем-то, «не фонтан», – хотя по Критерию №1 он очень даже «гут».)  = «Критерий №2».
Срединность же узла здесь даже не квалификационный, а классификационный критерий


\subsection{Операции перобразования узлов, введённые автром}

\begin{enumerate}
\item СПРЯМЛЕНИЕ  – процедура, превращающая Травяной узел в Прямой. 
\item РЕВЕРСИРОВАНИЕ  – технически, это не узловязочная операция, это метод создания нового узла, и заключается он в том, что ходовой и коренной концы верёвки (одной или обеих) меняются "должностями".
\item ИНВЕРТИРОВАНИЕ (ходовых) или ИНВЕРСИЯ  – узловязочная операция, при которой ходовые концы на крайнем их пересечении меняются местами. Пример: взаимо-инвертируемы узлы Травяной и Тёщин 
\item СВИВКА - Из «хантера» при помощи свивки мы получаем Инфинити
\item операция, обратная спрямлению (автор названия не дал)
\end{enumerate}


\paragraph{Спрямление}
\begin{itemize}
\item травяной - прямой
\item трасформатор - хирургический
\end{itemize}

Спрямление лучше определять через обратную операцию. Операция, обратная спрямлению, определена для узлов, представляющих собой комбинации полуузлов: прямого, бабьего, хирургического, академического - и заключается в заведении ходового конца между компонентами нижнего полуузла. Верхний полуузел при этом распадается.
Соответственно, спрямление заключается в извлечении ходового конца из окна в нижнем полуузле, и проведении его вместе с коренным. Если спрямление выполняется с обоими ходовыми концами, то они должны образовать между собой новый полуузел.

\paragraph{Инверсия}
\begin{itemize}
\item Травяной - Тёщин
\item Узел Эшли(\#1452) - Скорняжный 
\item Альтер-С(\#1408) - \#1409
\item Хантер - Колобок Голдобина (сомнительно)
\item Одновитковый змеиный - Пропеллер (\#1421)
\end{itemize}


\paragraph{Инфинити - Хантер - Колобок Голдобина}
Автор утверждает, что Хантер и колобок Голдобина связаны инверсией, а Хантер и Инфинити - свивкой.
В действительности, легко видеть, что операция, превращающая Хантер в колобок Голдобина, отличается от операции, превращающей Хантер в Инфинити, только направлением создаваемых пересечений. А операция, превращающая узел Эшли(\#1452) в Скорняжный, на них не похожа.

Узлы левые - скручивание левое - колобок\\
Узлы левые - скручивание нулевое - Хантер\\
Узлы левые - скручивание правое - Инфинити

Узлы правые - скручивание правое - колобок



\includegraphics[scale=0.5]{images/hunter-bend.eps}
\includegraphics[scale=0.5]{images/kolobok-goldobina.eps}
\includegraphics[scale=0.5]{images/infinity.eps}


\subsection{Система именования узлов и их элементов, введённая автором}
Для обозначения узлов и их элементов автор использует систему сокращений:

\begin{tabular}{ll}
	ХК & Ходовой конец \\
	КК & Коренной конец \\
	ОО & Открытая опора \\
	ЗО & Закрытая опора \\
	КУ & Контрольный узел \\
	ПШ & Полуштык \\
	ИП & Исходное положение
\end{tabular}

\begin{tabular}{ll}
	СП & Срединная петля \\
	НП & Незатягивающаяся петля \\
	2НП & Двойная незатягивающаяся петля \\
	РП & Регулируемая петля \\
	ПЗ & Петля-зящёлка \\
	ПУ & Петля-удавка \\
	хич & Крепёжный узел \\
	РХ & Узел для кольца (ринг-хич) \\
	СУ & Схватывающий узел \\
	РСУ & Реверсно-сбрасываемый узел (самосброс) \\
	СтУ & Стопорный узел
\end{tabular}
